\chapter{Joint Swelling}

\section*{Definition}
Joint swelling (joint effusion) is the abnormal accumulation of fluid within or around a joint capsule, presenting as visible or palpable enlargement. It may be inflammatory, non-inflammatory, or hemorrhagic, and accurate classification requires arthrocentesis and fluid analysis.

\section*{What Happens in Your Body}
Synovial fluid is produced by the synovial membrane as an ultrafiltrate of plasma supplemented with hyaluronic acid and lubricin. Inflammation increases synovial capillary permeability, allowing plasma proteins and leukocytes to enter the joint space. In non-inflammatory conditions, the fluid is transudative with low cell counts. Hemarthrosis results from direct trauma or bleeding diathesis, introducing blood and breakdown products into the joint.

\section*{Causes}
\begin{itemize}
  \item \textbf{Osteoarthritis:} Non-inflammatory effusion from cartilage debris and mechanical irritation
  \item \textbf{Rheumatoid arthritis:} Symmetric inflammatory effusion with synovial hypertrophy
  \item \textbf{Crystal arthropathy:} Neutrophil-rich effusion in gout (urate crystals) and pseudogout (calcium pyrophosphate crystals)
  \item \textbf{Septic arthritis:} Purulent effusion requiring urgent drainage and antibiotics
  \item \textbf{Traumatic:} Hemarthrosis from ligament tear, fracture, or meniscal injury
  \item \textbf{Hemorrhagic:} Coagulopathy (hemophilia, von Willebrand disease), anticoagulant therapy
  \item \textbf{Reactive:} Post-infectious (Reactive arthritis, rheumatic fever), enteropathic arthritis
  \item \textbf{Neoplastic:} Pigmented villonodular synovitis (PVNS), synovial chondromatosis
\end{itemize}

\section*{When to See a Doctor}
See your doctor if:
\begin{itemize}
  \item Acute monoarticular swelling with erythema, warmth, and fever (rule out septic arthritis)
  \item Joint swelling after trauma with inability to bear weight
  \item Chronic swelling with morning stiffness lasting > 30 minutes
  \item Rapid recurrence of effusion after aspiration
\end{itemize}

\section*{Related Symptoms}
\begin{itemize}
  \item Joint pain and tenderness
  \item Limited range of motion due to capsular distension
  \item Warmth and erythema over the joint
  \item Ballotable patella sign in knee effusion
  \item Crepitus, instability, or locking in mechanical joint derangement
\end{itemize}
\textbf{Important:} Warm, erythematous, monoarticular swelling is septic arthritis until proven otherwise; urgent arthrocentesis is mandatory.

\section*{Self-Care / Home Management}
\begin{itemize}
  \item Rest, ice, compression, and elevation (RICE) for acute traumatic swelling
  \item Avoid weight-bearing on the affected joint until evaluated
  \item Over-the-counter NSAIDs for pain and inflammation
  \item Elastic bandage or compression wrap to limit further fluid accumulation
\end{itemize}

\section*{How Your Doctor Will Evaluate You}
\begin{itemize}
  \item Arthrocentesis with synovial fluid analysis: cell count, Gram stain, culture, crystal analysis
  \item Serum laboratory studies: CBC, ESR, CRP, RF, anti-CCP, uric acid, ANA
  \item Plain radiography to assess joint space, erosions, osteophytes, and fracture
  \item Ultrasound to confirm effusion, guide aspiration, and assess synovial hypertrophy
  \item MRI for detailed evaluation of soft tissue, menisci, ligaments, and occult fracture
\end{itemize}

\section*{Treatment Options}
\begin{itemize}
  \item Antibiotics and surgical drainage for septic arthritis
  \item NSAIDs, colchicine, or corticosteroids for crystal-induced arthritis
  \item NSAIDs, DMARDs, or biologics for inflammatory arthritis
  \item Corticosteroid injection after infection is excluded
  \item Arthrocentesis for symptomatic relief of tense effusions
  \item Surgical intervention for structural lesions (meniscal tear, ligament repair, synovectomy)
\end{itemize}

\section*{References}
\begin{thebibliography}{99}
\bibitem{ref1} Brannan SR, Jerrard DA. "The evaluation and treatment of acute monoarticular arthritis." \textit{Emergency Medicine Clinics of North America}, 2008;26(3):735--753.
\bibitem{ref2} Margaretten ME, Kohlwes J, Moore D, et al. "Does this adult patient have septic arthritis?" \textit{JAMA}, 2007;297(13):1478--1488.
\bibitem{ref3} Courtney P, Doherty M. "Joint aspiration and injection and synovial fluid analysis." \textit{Best Practice \& Research Clinical Rheumatology}, 2013;27(2):137--149.
\bibitem{ref4} Shmerling RH. "Synovial fluid analysis: a critical appraisal." \textit{Rheumatic Disease Clinics of North America}, 1994;20(2):503--512.
\bibitem{ref5} Punzi L, Oliviero F. "Arthrocentesis and synovial fluid analysis in the diagnosis and management of arthritis." \textit{Nature Reviews Rheumatology}, 2011;7(1):43--52.
\bibitem{ref6} Roberts WN, Hayes CW, Breitbach SA, et al. "Dry taps and what to do about them: a pictorial essay." \textit{Journal of Rheumatology}, 1996;23(11):1971--1975.

\end{thebibliography}
